\pageheading{7 Preserve the semantic guards that make the proofs true}
\textbf{Entry into halt is still an active step.} A live transition with next state \code{none} writes at the old head and moves once. At the destination it reads the \emph{original destination bit}, because the write occurred at the distinct old head. Receive target 215 is valid and yields south color 432 or 433; it must not be rejected as an out-of-range operating state. The maintained statement is \src{WangMachineStep}{115}{123}, and both proposed row interfaces retain these cases.

\textbf{An unchanged plain cell can locally finish.} The tuple $\langle0,0,0,438\rangle$ is a genuine finish alternative with plain north and zero horizontal edges. Those conditions alone therefore do not imply pass. Receiver causality deliberately yields pass \emph{or} finish. The stronger pass conclusion uses \code{no_south438}, whose proof depends on temporal matching and the live head. A finite local row containing a head, its halt receiver and a separate finish illustrates the forbidden local strengthening; it is not claimed to satisfy every persistent-strip or temporal premise.

\textbf{Do not strengthen observation to configuration equality.} In the interval $(-2,3)$, configurations with state \code{some 1}, head 1, and tapes \code{[0]} versus \code{[0,100]} have identical observed colors but different tapes. Even equality of all integer reads does not imply equality of the concrete tape lists: \code{[0]} and \code{[0,0]} are observationally equal and list-unequal. The complete-tile proposal proves tuple equality only at observed interior positions. It cannot recover exterior tape content, a canonical tape representation or tile-index injectivity.

\textbf{The initializer is not a blank configuration.} The actual clock-one configuration has tape \code{[0]}, head 1 and state \code{some 1}; its colors are 1 at zero, 4 at one, and 0 elsewhere. See \src{DecimalMachineInitialization}{20}{52} and \src{WangMachineEncoding}{147}{169}. The head initializer is the literal $\langle3,436,4,4\rangle$, not a similar tuple with south 188, as specified in \src{WangTiles942InitialStrip}{23}{34}. No proposal replaces this with an arbitrary initial state or independently selected witness.

\textbf{Keep spatial and temporal dependencies separate.} Receiver location needs allowed tiles, horizontal matching, persistent walls, the original encoded row and its live state. It does not need temporal matching. Complete tile equality does need the temporal information that rules out an unchanged finish. Moving both results into one stronger context would hide this useful distinction and make the causal theorem unnecessarily difficult to reuse.

All positions remain signed integers. The actual destination is the one-cell move of the old head, not a natural-index subtraction. Wall-edge bounds retain their asymmetric values, and no statement extends the interior equality to seam walls or exterior cells. The original destination tape bit remains unconstrained; no blank-tape premise is added to a general row theorem.

\textbf{Respect the newer endpoint without overstating it.} The late maintained source already derives marked-language consistency equivalence and the forward consistency-to-polynomial-unsolvability implication. The retained local proposals do not reprove those results and are not prerequisites for them. They also do not provide the remaining reverse arithmetic/controller construction or an internal-ZFC formalization. Their scope is proof maintenance and stronger actual-row interfaces only.
